\subsection{The 't Hooft - Polyakov Monopole}

Let us look for time-independent, finite energy solutions in the
Georgi-Glashow model. Finiteness of energy requires that as
$r\rightarrow \infty$, the energy density $\theta_{00}$ given by
(\ref{XXX}) must approach zero faster than $1/r^3$. This means that
as $r\rightarrow \infty$, our solution must go over to a Higgs vacuum
defined by (\ref{XXXI}).  In the following, we will first assume that
such a finite energy solution exists and show that it can have a
monopole charge related to its soliton number which is, in turn,
determined by the associated Higgs vacuum. This result is proven
without having to deal with any particular solution explicitly. Next,
we will describe the 't Hooft-Polyakov ansatz for explicitly
constructing one such monopole solution. We will also comment on the
existence of Dyonic solutions. For convenience, in this section we
will use the vector notation for the $SO(3)$ gauge group indices and
not for the spatial indices. 

\noindent\underline{The Topological Nature of Magnetic Charge:}
\noindent Let $\vec{\phi}_{vac}$ denote the field $\vec{\phi}$ in a
Higgs vacuum. It then satisfies the equations
\begin{eqnarray}
&\vec{\phi}_{vac}\cdot\vec{\phi}_{vac} = a^2\,, &\nonumber \\
&\partial_{\mu} \vec{\phi}_{vac}-e\,\vec{W}_{\mu} \times
\vec{\phi}_{vac} = 0\,,& 
\label{vec-vac}
\end{eqnarray}
which can be solved for $\vec{W}_\mu$. The most general solution is
given by 
\begin{equation}
\vec{W}_{\mu} = \frac{1}{ea^{2}}\, \vec{\phi}_{vac} \times
\partial_\mu\vec{\phi}_{vac}+\frac{1}{a}\vec{\phi}_{vac}A_{\mu}\,.
\label{sol-vac}
\end{equation}
To see that this actually solves (\ref{vec-vac}), note that 
$\partial_\mu\vec{\phi}_{vac}\cdot\vec{\phi}_{vac}=0$, so that
$$
\frac{1}{ea^2}(\vec{\phi}_{vac}\times\partial_{\mu}\vec{\phi}_{vac})
\times \vec{\phi}_{vac}= \frac{1}{ea^2}\left(\partial_\mu\vec{\phi}_{vac}a^2
-\vec{\phi}_{vac} (\vec{\phi}_{vac}\cdot\partial_\mu\phi_{vac})\right)= 
\frac{1}{e}\partial_\mu\vec{\phi}_{vac}\,.
$$
The first term on the right-hand side of Eq. (\ref{sol-vac}) is the
particular solution, and $\vec{\phi}_{vac} A_\mu$ is the general
solution to the homogeneous equation. Using this solution, we can now
compute the field strength tensor $\vec{G}_{\mu\nu}$. The field
strength $F_{\mu\nu}$ corresponding to the unbroken part of the gauge
group can be identified as  
\begin{eqnarray}
F_{\mu\nu}&=&\frac{1}{a}\vec{\phi}_{vac}\cdot\vec{G}_{\mu \nu}
\nonumber\\
&=&\partial_\mu A_\nu -\partial_\nu A_\mu +
\frac{1}{a^3 e}\vec{\phi}_{vac}\cdot (\partial_\mu\vec{\phi}_{vac}
\times\partial_\nu\vec{\phi}_{vac})\,.
\label{F-vac}
\end{eqnarray}
Using the equations of motion in the Higgs vacuum it follows that  
$$
\partial_\mu F^{\mu\nu}=0\,,\qquad \partial_\mu\,\widetilde F^{\mu\nu}=0\,.
$$
This confirms that $F_{\mu\nu}$ is a valid $U(1)$ field strength
tensor. The magnetic field is given by
$B^i=-\frac{1}{2}\epsilon^{ijk}F_{jk}$. Let us now consider a static,
finite energy solution and a surface $\Sigma$ enclosing the core of
the solution. We take $\Sigma$ to be far enough so that, on it, the
solution is already in the Higgs vacuum. We can now use the magnetic
field in the Higgs vacuum to calculate the magnetic charge $g_\Sigma$
associated with our solution: 
\begin{equation}
g_\Sigma=\int_\Sigma\,B^i ds^i=-\frac{1}{2 ea^3} \int_\Sigma\, 
\epsilon_{ijk}\,\vec{\phi}_{vac}\cdot \left(\partial^j\vec{\phi}_{vac}
\times\partial^k\vec{\phi}_{vac}\right)ds^i\,.
\label{XL}
\end{equation}
It turns out that the expression on the right hand side is a
topological quantity as we explain below: Since $\phi^2=a$; the
manifold of Higgs vacua (${\cal M}_0$) has the topology of $S^2$. 
The field $\vec{\phi}_{vac}$ defines a map from $\Sigma$ into 
${\cal M}_0$. Since $\Sigma$ is also an $S^2$, the map 
$\phi_{vac}:\Sigma\rightarrow{\cal M}_0$ is characterized by its
homotopy group $\pi_2 (S^2)$. In other words, $\phi_{vac}$ is
characterized by an integer $\nu$ (the winding number) which counts the
number of times it wraps $\Sigma$ around ${\cal M}_0$. In terms of
the map $\phi_{vac}$, this integer is given by
\begin{equation}
\nu = \frac{1}{4\pi a^{3}} \, \int_\Sigma\,\frac{1}{2}
\epsilon_{ijk}\vec{\phi}_{vac}\cdot\left(\partial^j\vec{\phi}_{vac}
\times \partial^k\vec{\phi}_{vac}\right)\,d s^i\,.
\label{XLI}
\end{equation}
Comparing this with the expression for magnetic charge, we get the
important result
\begin{equation}
g_{\Sigma} = - \frac{4\pi \nu}{e}\,.
\label{XLII}
\end{equation}
Hence, the winding number of the soliton determines its monopole
charge. Note that the above equation differs from the Dirac
quantization condition by a factor of $2$. This is because the
smallest electric charge which could exist in our model is 
$q_0=e\hbar/2$ in terms of which, (\ref{XLII}) reduces to the Dirac
condition. 

\noindent\underline{An Ansatz for Monopoles:}
\noindent Now we describe an ansatz proposed by 't Hooft \cite{thooft}
and Polyakov \cite{polya} for constructing a monopole solution in the 
Georgi-Glashow model. For a spherically symmetric, parity-invariant,
static solution of finite energy, they proposed:
\begin{eqnarray}
\phi^a & = & \frac{x^a}{er^2}\,H(aer)\,,\nonumber\\
W^a_i&=&-\epsilon^a_{ij}\frac{x^j}{er^2}\left(1 - K(aer)\right)\,,
\qquad W^a_0= 0\,.
\label{tp-ansatz}
\end{eqnarray}
For the non-trivial Higgs vacuum at $r\rightarrow\infty$, they chose
$\phi^c_{vac}=ax^c/r= a\hat{x}^c$. Note that this maps an $S^2$ at
spatial infinity onto the vacuum manifold with a unit winding number.
The asymptotic behaviour of the functions $H(aer)$ and $K(aer)$ are
determined by the Higgs vacuum as $r\rightarrow\infty$ and regularity
at $r=0$. Explicitly, defining $\xi=aer$, we have: as
$\xi\rightarrow\infty,\,H\sim\xi,\, K\rightarrow 0$ and as
$\xi\rightarrow 0,\, H\sim\xi,\,(K-1)\sim\xi$. The mass of this
solution can be parametrized as
$$
M = \frac{4 \pi a}{e} f \, (\lambda /e^{2})
$$
For this ansatz, the equations of motion reduce to two coupled
equations for $K$ and $H$ which have been solved exactly only in
certain limits.  
For $r\rightarrow 0$, one gets $H\rightarrow ec_1 r^2$ and
$K=1+ec_2r^2$ which shows that the fields are non-singular at
$r=0$. For $r\rightarrow\infty$, we get $H\rightarrow\xi+c_3
exp(-a\sqrt{2\lambda}r)$ and $K\rightarrow c_4\xi exp(-\xi)$ which
leads to $W^a_i\approx-\epsilon^a_{ij}x^j/er^2$. Once again, defining
$F_{ij}=\phi^c G^c_{ij}/a$, the magnetic field turns out to be  
$B^{i} = - x^i/e r^3$. The associated monopole charge is $g=-4\pi/e$,
as expected from the unit winding number of the solution. It should be
mentioned that 't Hooft's definition of the Abelian field strength
tensor is slightly different but, at large distances, it reduces to
the form given above. 

Note that in the above monopole solution, the presence of the Dirac
string is not obvious. To extract the Dirac string, we have to perform
a singular gauge transformation on this solution which rotates the
non-trivial Higgs vacuum $\phi^c_{vac}=a\hat{x}^c$ into the trivial
vacuum $\phi^c_{vac}= a \delta^{c3}$. In the process,the gauge field
develops a Dirac string singularity which now serves as the source of
the magnetic charge \cite{thooft}.  

\noindent\underline{The Julia-Zee Dyons:}

The 't Hooft-Polyakov monopole carries one unit of magnetic charge and
no electric charge. The Georgi-Glashow model also admits solutions
which carry both magnetic as well as electric charges. An ansatz for
constructing such a solution was proposed by Julia and Zee
\cite{Julia-Zee}. In this ansatz, $\phi^a$ and $W^a_i$ have exactly
the same form as in the 't Hooft-Polyakov ansatz, but $W^a_0$ is no
longer zero: $W^a_0=x^a J(aer)/er^2$. This serves as the source for
the electric charge of the dyon.  It turns out that the dyon electric
charge depends of a continuous parameter and, at the classical level,
does not satisfy the quantization condition. However, semiclassical
arguments \cite{TW1,GSW} show that, in CP invariant theories, and at
the quantum level, the dyon electric charge is quantized as $q=n\hbar
e$. This can be easiy understood if we recognize that a monopole is
not invariant under a guage transformation which is, of course, a
symmetry of the equations of motion. To treate the associated
zero-mode properly, the gauge degree of freedom should be regarded as
a collective coordinate. Upon quantization, this collective coordinate
leads to the existence of electrically charged states for the monopole
with discrete charges.  In the presence of a CP violating term in the
Lagrangian, the situation is more subtle as we will discuss later. In
the next subsection, we describe a limit in which the equations of
motion can be solved exactly for the 'tHooft-Polyakov and the
Julia-Zee ansatz. This is the limit in which the soliton mass
saturates the Bogomol'nyi bound.

\subsection{The Bogomol'nyi Bound and the BPS States}

In this subsection, we derive the Bogomol'nyi bound \cite{bogomol} on
the mass of a dyon in term of its electric and magnetic charges which
are the sources for $F^{\mu\nu}=\vec{\phi}\cdot \vec{G}^{\mu\nu}/a$. 
Using the Bianchi identity (\ref{bianchi}) and the first equation in
(\ref{XXIX}), we can write the charges as 
\begin{eqnarray} 
g&\equiv&\int_{S^2_\infty}B_i dS^i=\frac{1}{a}\int{\cal B}^a_i\phi^a
dS^i = \frac{1}{a} \int {\cal B}^a_i (D^i \phi)^a d^3x\,,
\nonumber\\
q&\equiv&\int_{S^2_\infty}E_i dS^i=\frac{1}{a}\int{\cal E}^a_i\phi^a
dS^i = \frac{1}{a} \int {\cal E}^a_i (D^i \phi)^a d^3x\,. 
\label{gq}
\end{eqnarray}
Now, in the center of mass frame, the dyon mass is given by 
$$
M\equiv \int d^3x \theta_{00}= \int d^3x \left(\,\frac{1}{2}
\left[({\cal E}^a_k)^2+({\cal B}^a_k)^2 + (D_k\phi^a)^2 + 
(D_0\phi^a)^2 \right]+ V(\phi) \right)\,,
$$
where, $\theta_{\mu\nu}$ is the energy momentum tensor. After a little
manipulation, and using the expressions for the electric and magnetic
charges given in (\ref{gq}), this can be written 
as 
\begin{eqnarray}
M&=&\int d^3x \left(\,\frac{1}{2}\left[({\cal E}^a_k -
D_k\phi^a\sin\theta)^2 + ({\cal B}^a_k-D_k\phi^a\cos\theta )^2
+(D_0\phi^a)^2\right] + V(\phi) \right)
\nonumber\\
&+& a (q \sin \theta + g \cos \theta )\,,
\label{M-bog}
\end{eqnarray}
where $\theta$ is an arbitrary angle.  Since the terms in the first
line are positive, we can write $M\geq (q \sin \theta + g \cos \theta
)$. This bound is maximum for $\tan\theta = q/g$. Thus we get the
Bogomol'nyi bound on the dyon mass as
$$
M \geq a\, \sqrt{g^2 + q^2}\,.
$$
For the 't Hooft-Polyakov solution, we have $q = 0$, and thus, $M \geq
a|g|$. But $|g|=4\pi/e$ and $M_W=ae\hbar=aq$, so that 
$$
M \geq a \frac{4 \pi}{e} = \frac{4 \pi}{e^{2}\hbar} M_{W} =
\frac{4 \pi \hbar}{q^{2}} M_{W} = \frac{\nu}{\alpha} M_{W}\,.
$$
Here, $\alpha$ is the fine structure constant and $\nu=1$ or $1/4$,
depending on whether the electron charge is $q$ or $q/2$. Since
$\alpha$ is a very small number ($\sim 1/137$ for electromagnetism),
the above relation implies that the monopole is much heavier than the
W-bosons associated with the symmetry breaking. 

From (\ref{M-bog}) it is clear that the bound is not saturated unless
$\lambda\rightarrow 0$, so that $V(\phi)=0$. This is the 
Bogomol'nyi-Prasad-Sommerfield (BPS) limit of the theory
\cite{bogomol,PS}. Note that in this limit,  $\phi_{vac}^2 =a^2$ is no
longer determined by the theory and, therefore, has to be imposed as
a boundary condition on the Higgs field. Moreover, in this limit, the
Higgs scalar becomes massless. Now, to saturate the bound we have to set 
\begin{equation}
D_0 \phi^a=0\,,\qquad {\cal E}^a_k=(D_k\phi)^a\sin\theta\,,
\qquad{\cal B}^a_k = (D_k\phi)^a\cos\theta\,,
\label{L}
\end{equation}
where, $\tan\theta=q/g$. In the BPS limit, one can use the 
't Hooft-Polyakov (or the Julia-Zee) ansatz either in (\ref{XXIX}), or
in (\ref{L}) to obtain the exact monopole (or dyon) solutions
\cite{bogomol, PS}. These solutions automatically saturate the
Bogomol'nyi bound and are referred to as the BPS states. Also, note
that in the BPS limit, all the perturbative excitations of the theory
saturate this bound and, therefore, belong to the BPS spectrum. As we
will see later, the BPS bound appears in a very natural way in
theories with $N=2$ supersymmetry. 

\subsection{Monopoles from a Distance}

Till now, we have described the monopoles arising in the
Georgi-Glashow model in terms of the structure of the  Higgs vacuum of
the theory. In this section, we will consider monopoles in a general
Yang-Mills-Higgs system and relate the Higgs vacuum description to a
description in terms of the unbroken gauge fields. These are the gauge
fields which remain massless and are relevant for the study of
monopoles at large distances. This formulation is convenient for
describing non-abelian monopoles. 

Let $\phi$ transform as a vector in a given representation of a gauge
group $G$. For convenience of notation, in the following we do not
distinguish between the group element $g$ and a given realization of
it. Writing the gauge fields as $W_\mu=T^aW^a_\mu$, we can construct
the covariant derivative of $\phi$ and the curvature tensor of $W_\mu$
as
\begin{eqnarray}
D_\mu \phi &=& \partial_{\mu} \phi + ie W_{\mu} \phi\,,
\nonumber \\
G_{\mu\nu}&=&\partial_\mu W_\nu -\partial_\nu W_\mu +
ie [W_\mu, W_\nu] \,.
\nonumber
\end{eqnarray}
The Lagrangian density ${\cal L}$, the stress-energy tensor  
$\theta_{\mu\nu}$ and the gauge current $j^a_\mu$ are then given by  
\begin{eqnarray}
{\cal L}&=&-\frac{1}{4}(G^a_{\mu\nu})^2+(D^\mu\phi)^\dagger D_\mu\phi
-V(\phi) \,,
\nonumber\\
\theta_{\mu\nu}&=&-G^a_{\mu\lambda}G^{a\lambda}_\nu +
(D_\mu\phi)^\dagger(D_\nu\phi) + (D_\nu\phi)^\dagger(D_\mu\phi) 
- g_{\mu\nu} {\cal L} \,,
\nonumber \\
\theta_{00}&=& \frac{1}{2}{\cal E}^{ia}{\cal E}^a_i +
\frac{1}{2}{\cal B}^{ia}{\cal B}^a_i +D_0\phi^\dagger D_0\phi +
D_i\phi^\dagger D_i\phi + V(\phi)\,,
\nonumber \\
j^a_\mu &=& ie \phi^\dagger T^a D_\mu\phi -
ie (D_\mu\phi)^\dagger T^a \phi\,.
\label{LIII}
\end{eqnarray}
Here, the Higgs potential is gauge invariant: $V(g\phi)=V(\phi)$. 
The equations of motion following from the above Lagrangian are 
\begin{equation}
D^\mu D_\mu\phi^a=-\frac{\partial V}{\partial \phi^a}\,,\qquad
D^\nu G^a_{\mu\nu} = - j_\mu^a\,.
\label{LIV}
\end{equation}
When the gauge group is $SO(3)$ spontaneously broken to $U(1)$, we can
work out the Bogomol'nyi bound exactly as in the previous section and
the outcome is 
\begin{equation}
M\geq \sqrt 2 \, a \sqrt{g^2 + q^2}\,.
\label{Bbound}
\end{equation}
For a general gauge group $G$, the Higgs vacuum, as in the
Georgi-Glashow model, is defined by  
$$
V(\phi)=0\,,\qquad D_\mu\phi=0 \,.
$$
The first equation defines the vacuum manifold ${\cal M}_0\equiv 
\{\phi\,: V(\phi)=0\}$, and the second equation leads to 
$$
[D_\mu, D_\nu] \phi =G_{\mu\nu} \phi = 0\,. 
$$
Thus, in the Higgs vacuum, $G_{\mu\nu}$ takes values in a subgroup of
the gauge group $G$ which keeps the Higgs vacuum invariant. We denote
this unbroken subgroup of $G$ by $H$. The generators of $G$ which do
not keep the Higgs vacuum invariant are of course broken and the
corresponding gauge bosons become massive.  If $V(\phi)$ does not have
extra global symmetries and accidental minima, then it is reasonable
to assume that the action of $G$ on ${\cal M}_0$ is transitive. This
means that any point in ${\cal M}_0$ is related to any other point
(and, in particular, to a reference point $\phi_0$) by some element of
$G$. Therefore, the little group or the invariance group, $H\subset
G$, of any point in ${\cal M}_0$ is isomorphic to the the little group
of any other point. Hence, the structure of ${\cal M}_0$ is described
by the right coset $G/H$.

For a solution  to have finite energy, at sufficiently large distances
from the core of the solution the field $\phi$ must take values in
the Higgs vacuum. Let $\Sigma$ be a 2-dimensional surface around the
core such that, on this surface, $\phi$ is already in ${\cal M}_0$. 
On this surface, $\phi$ describes a map from $\Sigma$ (with the
topology of $S^2$) into ${\cal M}_0$. This map is characterised by its
homotopy class which has to be an element of $\pi_2({\cal M}_0)\simeq
\pi_2(G/H)$. As described before, the associated topological number is
the magnetic charge of the solution. As long as no monopole crosses
the surface $\Sigma$, $\phi$ remains a continuous function of time and its
homotopy class does 
not change. To show that the map $\phi$ satisfies the group properties
of $\pi_2(G/H)$, one has to consider several widely separated
monopoles and study how their magnetic charges combine. For a
discussion of this issue, see \cite{coleman1}. 

The above discussion of the topological characterisation of the
monopole is in terms of the structure of the Higgs vacuum. However, it
is more natural to have a description in terms of the unbroken gauge
fields. The relationship  between these two descriptions is contained
in the equation $D_\mu\phi =0$ which is valid on the $S^2$ surface
$\Sigma$. Let us parametrise $S^2$ by a square $\{ 0 \leq s, t \leq 1
\}$. The map $r(s,t)$ from this square to the sphare is single valued
everywhere except on the boundary of the square which is identified
with a single point on the sphere: 
$$
r(0,t) = r(1,t) =r(s,0) =r(s,1) = r_0\,.
$$
For fixed $s$, as $t$ varies from $0$ to $1$, $r(s,t)$ describes a
closed path on $S^2$. 

% ==== FIGURE_SLOT_2 ====
The covariant derivative along $t$ becomes
$$
D_t \phi(s,t) =\frac{\partial r^i}{\partial t}D_i \phi(s,t)=0\,.
$$
For constant $s$, this can be solved as 
$$
\phi(s,t)=g(s,t)\phi(s,0)\,, 
$$
where, $g(s,t)$ satisfies $D_t g=0$, $g(s,0)=1$ and is given by the
path-ordered integral 
\begin{equation}
g (s,t) =  P\left(\,\exp \left(\,ie\int^{t}_{0}
\vec{W}\cdot\frac{\partial \vec{r}}{\partial t} \, dt\right)\right)\,. 
\label{LVIII}
\end{equation}
Clearly, $g(s,t)$ is an element of $G$ which gauge transforms any
$\phi(s,t)$ into a reference $\phi(s,0)$ and  $\phi(s,0)=\phi_0$ is
independent of $s$. Therefore, $g(s,t)$ contains the same topological
information as $\phi(s,t)$ and is characterised by a homotopy class
in $\pi_2(G/H)$. Since at $s=0,1$ we have $\partial r/ \partial t =0$,
then, $g(0,t)=g(1,t)=1$. For $t=1$, $\phi(s,1)=g(s,1)\phi_0$. But
since $\phi(s,1)=\phi_0$, we conclude that $g(s,1)$ must be an element
of the unbroken gauge group $H$ and denote it by $h(s)$:
\begin{equation}
h(s)\equiv g(s,1)=P\,\exp\left(ie\int^1_0\,\vec{W}\cdot\frac{\partial
\vec{r}}{\partial t}\,dt\right)\,. 
\label{LIX}
\end{equation}
Since $h(1)=h(0)=1$, as we move along $s$, $h(s)$ describes a closed
path in $H$ and is thus, an element of $\pi_1(H)$. However, note that
as we vary $t$ from $1$ to $0$, $g(s,t)$ continuously interpolates
between $h(s)$ and the identity. Therefore, $h(s)$ can describe only
those closed paths in $H$ which are homotopic to the trivial path when
$H$ is embedded in $G$. We denote the subgroup of $\pi_1(H)$ which 
corresponds to such paths by $\pi_1(H)_G$. If the homotopy class of
the map $\phi(s,t)$ is changed in $\pi_2(G/H)$, the homotopy class of
$h(s)$ changes in $\pi_1(H)_G$. In fact, there is a one-to-one
correspondence between $\pi_2(G/H)$ and $\pi_1(H)_G$ and the two
groups are isomorphic \cite{coleman1}. Thus, we can equally well
characterise the monopole by its homotopy class in $\pi_1(H)_G$.   

Let us discuss this in some more detail. Since any closed path in $H$
is also a closed path in $G$, there is a 
natural homomorphism from $\pi_1(H)$ into $\pi_1(G)$. As the
discussion above shows, $\pi_1(H)_G$ is in the kernel of this
homomorphism. Moreover, for a compact connected group $G$,
$\pi_2(G)=0$, which implies that $\pi_2(G)$ can be embedded in
$\pi_2(G/H)$ as its identity element. What we have done above,
basically, is to construct part of the following exact sequence:
$$
\pi_2(G)\rightarrow\pi_2(G/H)\rightarrow\pi_1(H)\rightarrow 
\pi_1(G)\rightarrow\pi_1(G/H)\rightarrow\pi_0(H)
$$
Note that if $G$ is simply connected then, $\pi_1(G)=0$ and 
$\pi_1(H)_G=\pi_1(H)$. So that the full $\pi_1(H)$ enters the physical
description of the monopole. For non-simply connected $G$, this
possibility can be realised in the presence of a Dirac string. Such a
string appears as a singular point on the $(s,t)$ plane, in the
presence of which, it is no longer possible to continuously deform 
$h(s)$ to the identity map. Therefore, the homotopy class of $h(s)$ is
no longer restricted to $\pi_1(H)_G$. However, 
for non-simply connected $G$, it is still possible to have a
description in terms of an unrestricted $\pi_1$ group provided we
embed $G$ in its universal covering group $\tilde G$. Let us denote
the little group of $\phi$ by $\hat H$. Since $\tilde G$ is
necessarily connected, $\pi_1(\tilde G)=0$ and $\pi_1(\hat H)= 
\pi_1(H)_G$. As an example, let us consider the Georgi-Glashow
model. Here, $G = SO(3)$ (with $(Ta)_{ij}=-i\epsilon_{aij}$) and  
$\tilde{G} = SU(2)$ (with $T_a=\frac{1}{2}\sigma_a$). The
homotopically distinct paths in $H$ (or $\hat H$) are: 
$$
h(t)=\exp(it\,\vec{\phi}\cdot\vec{T}\,4\pi N/a),\quad {\rm for}\,\,\, 
0 \leq t \leq 1\,.
$$
The integer $N$, which  characterises the elements of $\pi_1$, is
determined as follows: 
$$
\begin{array}{llll}
\tilde G=SU(2),& \tilde H=U(1),& \vec{\phi}\cdot\vec{T}\in 
\frac{1}{2}{\bf Z},& \Rightarrow\,\, N \in {\bf Z}\,, \cr
G=SO(3),& H=SO(2),& \vec{\phi}\cdot\vec{T}\in {\bf Z},&
\Rightarrow\,\, N \in \frac{1}{2}{\bf Z}\,.
\end{array}
$$
Since $SO(3)\sim S^2/{\bf Z}_2$, only paths with integer $N$
contribute in both cases. Later we show that $g=4\pi N/e$. For $SU(2)$,
$q_0=e\hbar/2$ and $gq_0/4\pi\hbar=N/2$. For $SO(3)$, $q_0=e\hbar$ and
$gq_0/4\pi\hbar=N$.

For Glashow-Weinberg-Salam model, $G=SU(2)\times U(1)$ and $H=U(1)$ is
a linear combination of $SU(2)$ and $U(1)$.  Although $\pi_{1} (H) =
{\bf Z}$, $\pi_{1} (H)_{G} = 0$. Therefore, in this model, any
non-trivial monopole must have a Dirac string.

We set out to describe the monopole in terms of the unbroken gauge
fields (the $H$-fields). Although, we have obtained a description in
terms of $\pi_1(H)_G$, it is not manifest that $h(s)$, as given by
(\ref{LIX}), involves only the $H$-fields. We show this in the
following: The quantity $g^{-1}(s,t)D_s g(s,t)$ is invariant under a
$t$-dependent gauge transformation. Moreover, by construction, $D_t
g(s,t)=0$. Hence, we can write
\begin{equation}
\partial_t(g^{-1}D_s g)= D_t(g^{-1}D_s g)=g^{-1}[D_t,D_s] g =
ie g^{-1} G_{ij} g \frac{\partial r^i}{\partial t} \frac{\partial r^j}
{\partial s}\,.
\label{LX}
\end{equation}
Let us integrate the first and the last terms above from $t=0$ to
$t=1$. Since $g^{-1}D_s g=0$ at $t=0$ and $g^{-1}D_s g=h^{-1}dh/ds$ at
$t=1$, we get
\begin{equation}
h^{-1} \, \frac{dh}{ds} = ie \, \int^{1}_{0} dt \,
g^{-1} G_{ij}\,g \,\frac{\partial r^i}{\partial t} \,
\frac{\partial r^j}{\partial s}\,.
\label{LXI}
\end{equation}
Since $G_{ij}$ was calculated on $\Sigma$, it involves only the
$H$-gauge fields. The conjugation by $g$ does not bring in a
dependence on the massive gauge fields as is evident from the
left-hand side of the equation. Hence the map is given entirely in
terms of the $H$-fields without any reference to the Higgs field. As a
simple application, consider the Dirac monopole. Integrating
(\ref{LXI}) from $s=0$ to $s=1$, we get $h(1)=exp(ie\int {\vec
  B}\cdot{\vec ds})$. Since $h(1)=1$, this leads to the Dirac
quantization condition $eg=2\pi n$ or $qg/4\pi\hbar=n/2$. Another
interesting consequence of (\ref{LXI}) is a possible explanation for
the fractional charges of quarks. We describe this in the next
section. 

\subsection{The Monopole and Fractional Charges}

So far, we have seen how the existence of a monopole can quantize the
electric charge in integer units. In the physical world, however, we
also come across fractionally charged quarks. In the following, we see
how the existence of a monopole can also account for these fractional
charges \cite{CO,GO}.

Let us represent our adjoint Higgs by $\phi=\phi^a T^a$, where $T^a$
are the fundamental representation matrices. Moreover, we only
consider $\phi$ on the surface $\Sigma$ as described in the previous
subsection. With $\phi$ in the Higgs vacuum, a generator $T^a$ belongs
to the unbroken subgroup $H$ of $G$ provided $[T^a,\phi]=0$. This
implies that $\phi$ itself is a generator of $H$ and commutes with its
other generators. Thus the Lie Algebra of $H$ is of the form
$L(H)=u(1)\oplus L(K)$ and we choose $L(K)$ to be orthogonal to
$u(1)$: ${\rm Tr}(\phi K^a)=0$ for $K^a\in K$.  Locally, $H$ has the
structure $U(1)\times K$, though, this is not necessarily the global
structure. We refer to $K$ as the colour group and identify the $U(1)$
as corresponding to electromagnetism. The gauge fields in $H$ can
be decomposed as $W^\mu=A^\mu\phi/a + X^\mu$, with ${\rm Tr}(\phi
X^\mu)=0$. Expanding the covariant derivative $\partial_\mu+ ie W_\mu$
we can identify the electric charge operator which couples to $A_\mu$
as $Q=(e\hbar/a)\phi$.

Since $h^{-1} dh/ds$ is a generator of $H$, we may write
\begin{equation}
h^{-1} \frac{dh}{ds}=ie\alpha (s)\,\frac{\phi_0}{a} +
i \beta_{a} K^{a} \,.
\label{hab}
\end{equation}
Using ${\rm Tr}(T_a T_b)=\delta_{ab}$, $\phi(r)=g\phi_0 g^{-1}$ and
equation (\ref{LXI}), we get 
\begin{eqnarray}
\alpha (s)&=&-\frac{i}{ae}{\rm Tr}\left(\phi_0 h^{-1}\frac{dh}{ds} 
\right) \nonumber \\
&=&\frac{1}{a} \int^{1}_{0}{\rm Tr} \left(\varphi (r) G_{ij}\right) 
\frac{\partial r^i}{\partial t}\frac{\partial r^j}{\partial s}\,dt\,.
\nonumber
\end{eqnarray}
Identifying  the electromagnetic field strength tensor as
$F^{\mu\nu}={\rm Tr}\,(\phi G^{\mu\nu}/a)$, we get $\alpha(s)=
d\Omega/ds$, where, $\Omega(s)$ is the magnetic flux in a solid angle
subtended at the origin by the path $0\leq t\leq 1$ at fixed $s$ on
$S^2$. Substituting this back in (\ref{hab}) and integrating from
$s=0$  to $s$, gives
$$
h(s)= k(s) e^{i Q \Omega(s) / \hbar}\,.
$$
Since $h(1)=1$ and $\Omega (1) = g$ (where $g$ is the total magnetic
charge inside $\Sigma$), the quantization condition is replaced by 
$$
e^{igQ/\hbar} = k(1)^{-1}=k \in K \,.
$$
The left-hand side is invariant under $K$, therefore, we can at most
have $k\in {\bf Z}(K)$, the center of $K$. If we take $K = SU(N)$,
then, $k=e^{2\pi in/N}$ with $n=1,2,\dots N$. If all values of $n$ are
allowed, then $U(1)\cap K\subset Z_{N}$. This corresponds to the fact
that globally, $H$ cannot be decomposed as $U(1)\times K$. Now, let 
$|s>$ be a colour singlet. Then $k|s> = |s>$ and $exp(igq_s/\hbar)=1$,
which is again the Dirac quantization condition. Thus, if $q_0$ and
$g_0$ are the units of electric and magnetic charges for colour
singlets, then $q_s=n_sq_0$ and $g=mg_0$ with $g_0q_0=2\pi\hbar$. 
The colour non-singlet states $|c>$ can be classified according to
their behaviour under the center of the colour group $K$:
$$
k|c> = e^{2\pi i \, t(c)/N} |c> =e^{igQ / \hbar} |c>\,,
$$
where, $t(c)$ is an integer mod $N$. For a minimal monopole $g_{0}$,
we obtain $g_0 q_c /\hbar=2\pi(m + t(c)/N) $, hence,
$$
q_c = q_0 (m + \frac{t(c)}{N})\,. 
$$
If we set $N=3$, as for QCD, and $m=0$, then $q_c=q_0/3, 2q_0/3,q_0$. 

\subsection{Non-Abelian Magnetic Charge and the Montonen-Olive
  Conjecture} 

In this subsection we first consider the generalization of charge
quantisation to non-abelian monopoles \cite{GNO}, and then describe
the electric-magnetic duality conjecture of Montonen and Olive 
\cite{MO}.

Goddard, Nuyts and Olive \cite{GNO} attempted to classify all
$H$-monopole configurations. To describe such a monopole, we consider
a static configuration and choose the gauge $\vec{r}\cdot\vec{W}^a =0$.
At large distances, it is reasonable to write the magnetic componenets
of the field strength as
$$
G_{ij} = \frac{1}{4 \pi r^2} \, \epsilon_{ijk} \, \hat{r}^k G(r)\,, 
$$
where $D_\mu G(r)=0$. Since $G(r)$ transforms in the adjoint
representation, we can write $G(r)=g(s,t)G_{0}g(s,t)^{-1}$. 
Substituting the above expression in (\ref{LXI}), and integrating over
$s$, we get
\begin{eqnarray}
{\rm ln}\left(h(s)\right) & = & \frac{ie}{4\pi r^2} G_0  
\int^s_0 ds\int^1_0 dt\,\epsilon_{ijk}\hat{r}^k \,
\frac{\partial r^i}{\partial t}\frac{\partial r^j}{\partial s}
\nonumber \\
& = & \frac{ie}{4\pi} G_0\,\Omega (s)\,.
\label{LXX}
\end{eqnarray}
Here, $\Omega (s)$ is the solid angle subtended at the origin by the
loop $0\leq t\leq 1,\,s={\rm const}$ on $S^2$. The elements in
$\pi_1(H)_G$ are, therefore, given by  
$$
h(s)= e^{\frac{ie}{4 \pi} G_0\Omega(s)}\,. 
$$
Since, $h(1)=1$, the above equation implies that
\begin{equation}
e^{ie G_0} = 1\,,
\label{LXXI} 
\end{equation}
which is the generalized charge quantization condition.
Clearly, $G_0$ is arbitrary upto a conjugation in the gauge
group. This freedom can be used to solve the charge quantization
condition as follows: Assume that $H$ is compact and connected and let
$T$ denote an Abelian subgroup of $H$ generated by its Cartan
subalgebra. Then, any element of $H$ is conjugate to at least one
element of $T$. Thus, it is always possible to find a frame in which
$G_0 = \beta^a T^a$. The coefficients $\beta^a$ are still not unique
as they transform under the Weyl group of $H$ which keeps this
parametrization of $G_0$ unchanged. Therefore, the equivalence classes
of $\beta^a$, related by the action of the Weyl group, are the gauge
invariant objects which characterise the non-abelian magnetic
charges. The $\beta^a$ are determined by the quantization
condition $exp(ie \beta^a T^a)=1$. To solve this, let $\omega^a$
denote a weight vector of $H$ in the given representation, and let
$\Lambda(H)$ denote the weight lattice. Then the quantization
condition implies that
$$
e \beta^a \omega^a \in 2\pi{\bf Z}\,,\qquad {\rm for\, all}\,\, 
\omega\in \Lambda(H)\,.
$$      
Note that the factors of $e$ and $2\pi$ are convention dependent.  
Thus, $e\beta$ lies on a lattice dual to $\Lambda(H)$: $e\beta \in
\Lambda^*(H)$. This dual lattice can by itself be regarded as the
weight lattice of a dual group $H^v$ which has $e\beta$'s as its
weight vectors (For details, see \cite{GNO}). Moreover, $(H^v)^v=H$.
$H$ is referred to as the electric group and $H^v$ as the magnetic
group. The magnetic charges are related to $H^v$ in the same way that
electric charges are related to $H$. A simple example of a dual pair
of groups is provided by $SO(3)$ and $SU(2)$. In this case, $G_0=\beta
T^3$. $T^3$ has integral eigenvalues for $SO(3)$ and half-integral
eigenvalues for $SU(2)$. The quantization condition $exp(ie\beta
T^3)=1$ gives:  
\begin{eqnarray}
{\rm for}\,\,\,H=SO(3)\,:\quad e\beta &=& 4\pi \frac{n}{2}=4\pi \times 
({\rm weight \,\,of\,\,} SU(2))\,, \nonumber\\ 
{\rm for}\,\,\,H=SU(2)\,:\quad e\beta &=& 4\pi n = 4\pi \times 
({\rm weight\,\,of\,\,}SO(3))\,. \nonumber
\end{eqnarray}
For a general $SU(N)$ group, the dual relation is given by 
$(SU(NM)/{\bf Z}_N)^v = (SU(NM)/{\bf Z}_M)$ and, in particular,
$SU(N)^v = SU(N)/Z_N$.

Now, we will briefly describe the Montonen-Olive conjecture \cite{MO}
which is based on the above results. This  conjecture states that a
gauge theory is characterized by $H \times H^{v}$, and that we have
two equivalent descriptions of the theory: One in terms of $H$-gauge
fields with normal charged particles in the perturbative spectrum and
another, in terms of $H^{v}$-gauge fields with monopoles in the
perturbative spectrum. The Noether currents (associated with electric
charges) also get interchanged with topological currents (associated
with magnetic charges). Hence, the coupling constant $q/\hbar$ of the
$H$-theory is replaced by $g/\hbar$ in the in the $H^v$-theory. Since,
$g\sim 1/e$, this conjecture relates a strongly coupled theory to a
weekly coupled one, and vice-versa. As a result of this, it is
not easy to either prove or disprove this conjecture. Montonen and
Olive provided some semiclassical evidence in favour of this
conjecture in the BPS limit of the Georgi-Glashow model. This model
contains a Higgs boson, a photon and two massive charged vector bosons
in its perturbative spectrum, and magnetic monopoles as solitonic
classical solutions. The unbroken gauge group is self dual,
$H=U(1)=H^v$, therefore, the dual theory has the same form as the
original one with the monoples as elementary states and the massive
gauge bosons as solitonic solutions. The gauge boson mass in the dual
theory (where it appears as a soliton) can be computed using the BPS
formula and turns out to have the right value. Moreover, the
long-range force between two monopoles as obtained by Manton
\cite{manton}, can be obtained by calculating the potential between
$W$-bosons in the dual theory and turns out to be the same. 

Since the Montonen-Olive duality is non-perturbative in nature, it
cannot be verified in a perturbative framework unless we have some
kind of control over the perturbative and non-perturbative aspects of
the theory. Such a control is provided by superysymmetry. In fact, in 
the $N=4$ super Yang-Mills theory, some very non-trivial predictions
of this duality were verified in \cite{sen}. In later parts, we will
consider in detail the analogue of the Montonen-Olive duality in $N=2$
supersymmetric gauge theories. A prerequisite for this, however, is
the introduction of the $\theta$-term in the Yang-Mills action which
affects the electric charges of dyons.    

\subsection{The $\theta$-Parameter and the Monopole Charge}

In this section we will show, following Witten \cite{Witten79}, that
in the presence of a $\theta$-term in the Lagrangian, the magnetic
charge of a particle always contributes to its electric charge.

As shown by Schwinger and Zwanziger, for two dyons of charges
$(q_1,g_1)$ and $(q_2,g_2)$, the quantization condition takes the form
\begin{equation}
q_1g_2 - q_2g_1 = 2 \pi n \hbar
\label{LXXVI}
\end{equation}
For an electric charge $q_0$ and a dyon $(q_n,g_n)$, this gives
$q_0g_n = 2\pi n\hbar$. Thus, the smallest magnetic charge the dyon
can have is $g_0=2\pi\hbar/q_0$. For two dyons of the same magnetic
charge $g_0$ and electric charges $q_1$ and $q_2$, the quantization
condition implies $q_1-q_2= nq_0$. Therefore, although the difference
of electric charges is quantized, the individual charges are still
arbitrary. This arbitrariness in the electric charge of dyons can be
fixed if the theory is CP invariant as follows: Under a CP
transformation $(q,g)\rightarrow (-q,g)$. If the theory is CP
invariant, the existence of a state $(q,g_0)$ necessarily leads to the
existence of $(-q,g_0)$. Applying the quantization condition to this
pair, we get $2q =q_0\times integer$. This implies that $q=nq_0$ or
$q=(n+\frac{1}{2})q_0$, though at a time we can either have dyons of
integral or half-odd integral charge, and not both together.

In the above argument, it was essential to assume CP invariance to
obtain integral or half-odd inegral values for the electric charges of
dyons. However, in the real world, CP invariance is violated and there
is no reason to expect that the electric charge should be quantized as
above. To study the effect of CP violation, we consider the
Georgi-Glashow model with an additional $\theta$-term which is the
source of CP violation:
\begin{equation}
{\cal L}=-\frac{1}{4}\,F^a_{\mu\nu}F^{a\mu\nu}+\frac{1}{2}
(D_{\mu} \vec{\phi})^{2} - \lambda (\phi^{2} - a^{2} )^{2}
+\frac{\theta e^2}{32\pi^2}\,F^a_{\mu\nu}\tilde F^{a\mu\nu}\,. 
\label{LXXVIII}
\end{equation}
Here, $\tilde F^{a\mu\nu}=\frac{1}{2}\epsilon^{\mu\nu\rho\sigma} 
F^a_{\rho\sigma}$ and the vector notation is used to represent indices
in the gauge space. The presence of the $\theta$-term does not affect
the equations of motion but changes the physics since the theory is no
longer CP invariant. We want to construct the electric charge operator
in this theory. The theory has an $SO(3)$ gauge symmetry but the
electric charge is associated with an unbroken $U(1)$ which keeps the
Higgs vacuum invariant. Hence, we define an operator $N$ which
implements a gauge rotation around the $\hat{\phi}$ direction with
gauge parameter $\Lambda^a = \phi^a/a$. These transformations
correspond to the electric charge. Under $N$, a vector $\vec v $ and
the gauge fields $\vec A_\mu$ transform as    
$$
\delta \vec{v} = \frac{1}{a} \, \vec{\phi} \times \vec{v}\,, \quad
\delta \vec{A}_{\mu} = \frac{1}{ea} \, D_{\mu} \vec{\phi}\,.
$$
Clearly, $\vec{\phi}$ is kept invariant. At large distances where
$|\phi|=a$, the operator $e^{2\pi i N}$ is a $2 \pi$-rotation about
$\hat{\phi}$ and therefore $\exp\,(2\pi iN)=1$. Elsewhere, the
rotation angle is $2\pi |\phi|/a$. However, by Gauss' law, if the
gauge transformation is $1$ at $\infty$, it leaves the physical states
invariant. Thus, it is only the large distance behaviour of the
transformation which matters and the eigenvalues of $N$ are quantized
in integer units. Now, we use Noether's formula to compute $N$: 
$$
N=\int\,d^3x\left(\frac{\delta{\cal L}}{\delta\partial_0 A^a_i}\;
\delta A_{i}^{a} + \frac{\delta {\cal L}}{\delta \partial_{0}
\phi^{a}} \; \delta \phi^{a} \right )\,.
$$
Since $\delta \vec\phi= 0$, only the gauge part (which also includes
the $\theta$-term) contributes:
\begin{eqnarray}
\frac{\delta}{\delta \partial_{0}A^{a}_{i}} \,
\left ( F^{a}_{\mu \nu} F^{a\mu\nu}\right) &=&
4F^{aoi} = -4 {\cal E}^{ai}\,, \nonumber\\
\frac{\delta}{\delta \partial_{0}A^{a}_{i}} \,
\left(\tilde{F}^a_{\mu\nu}F^{a\mu\nu}\right) &=&
2\epsilon^{ijk}F^a_{jk} =-4 {\cal B}^{ai}\,. \nonumber
\end{eqnarray}
Thus, we get
\begin{eqnarray}
N &=&\frac{1}{ae} \int d^{3}x D_{i}\vec{\phi} \cdot \vec{\cal E}^i
- \frac{\theta e}{8\pi^{2}a}\,\int d^3x \, D_{i}\vec{\phi}
\cdot \vec{\cal B}^i \nonumber\\
&=& \frac{1}{e} Q - \frac{\theta e}{8 \pi^2} M \,,\nonumber
\end{eqnarray}
where, we have used equations (\ref{gq}). Here, $Q$ and $M$ are the
electric and magnetic charge operators with eigenvalues $q$ and $g$,
respectivly, and $N$ is quantized in integer units. This leads to the
following formula for the electric charge 
$$
q=ne + \frac{\theta e^2}{8\pi^2} g \,.
$$  
For the 't Hooft-Polyakov monopole, $n=1$, $g=-4\pi/e$, and
therefore, $q=e(1-\theta/2\pi)$. For a general dyonic solution we get
\begin{equation}
g= \frac{4\pi}{e} m,\qquad q = ne + \frac{\theta e}{2\pi}m\,.
\label{gqtheta}
\end{equation}
Thus, in the presence of a $\theta$-term, a magnetic monopole always
carries an electric charge which is not an integral multiple of some
basic unit.  

It is very useful to represent the charged states as points on the
complex plane, with electric charges along the real axis and magnetic
charges along the imaginary axis. A state can thus be represented as
\begin{equation}
q+ig=e(n+m\tau)\,,
\label{qig}
\end{equation}
where,
\begin{equation}
\tau=\frac{\theta}{2\pi} + \frac{4\pi i}{e^2}
\label{tau}
\end{equation}
In this parametrisation, the Bogomol'nyi bound (\ref{Bbound}) takes
the form 
\begin{equation}
M \geq {\sqrt 2}|ae(n+m\tau)|\,.
\label{BBtau}
\end{equation}
Note that (\ref{qig}) implies that all states lie on a two-dimensional
lattice with lattice parameter $\tau$ and (\ref{BBtau}) implies that
the BPS bound for a state is proportional to the distance of its
lattice point from the origin. These equations play a very important
role in the subsequent discussions.    
